\section{Results}

\subsection{Diffusivity spectra}

Figure~\ref{fig:spectra} shows the cohort MAP and NUTS spectra by tissue type and zone, with per-ROI posteriors in the supplementary atlas (Figure~\ref{fig:atlas}). The two estimators agreed closely throughout the cohort, with per-component cross-ROI correlations of $r = 0.99$, $0.94$, and $0.94$ on the three well-identified bins ($D = 0.25$, $3.0$, and $20\;\umsmm$) and $r = 0.51$--$0.73$ on the collinear intermediate bins.
From normal to tumor, both zones gained mass at $D = 0.25\;\umsmm$ at the expense of $D = 3.0\;\umsmm$ (and, in the transition zone, $D = 2.0\;\umsmm$ as well).
The normal baseline itself differed by zone. Peripheral-zone normal tissue was dominated by $D = 3.0\;\umsmm$ (mean $0.48$) with a wide cohort spread, whereas transition-zone normal tissue placed more, and less variable, mass on the $D = 2.0$ and $3.0\;\umsmm$ bins (each $\approx 0.29$).
Tumor spectra were likewise zone-dependent. The $D = 0.25\;\umsmm$ fraction was both higher and far more variable across the cohort in transition-zone tumor than in transition-zone normal ($0.25 \pm 0.14$ vs.\ $0.09 \pm 0.03$), while peripheral-zone tumor retained more $D = 2.0$--$3.0\;\umsmm$ mass than transition-zone tumor.
In the separate three-direction subcohort, per-direction spectra agreed closely at the restricted bin that carries the detection signal (mean coefficient of variation across directions $9\%$ at $D = 0.25\;\umsmm$, 13 ROIs) and diverged on the poorly identified bins, supporting trace averaging (Figure~\ref{fig:directions}).

\subsection{Tumor detection}

Table~\ref{tab:auc} and Figure~\ref{fig:roc} summarize detection performance.
ADC reached an AUC of $0.951$ (PZ) and $0.979$ (TZ) as a raw ranking, and $0.940$ and $0.964$ through the leave-one-out pipeline that also scores the spectral models (Methods).
The spectral classifiers matched it. Through that identical pipeline the eight-fraction logistic regression reached PZ $0.926$ / TZ $0.923$ (NUTS) and PZ $0.917$ / TZ $0.952$ (MAP), with bootstrap confidence intervals overlapping ADC's throughout.
Restricting the classifier to the restricted and free-water fractions $\{D = 0.25,\,3.0\;\umsmm\}$ did not reduce performance (NUTS PZ $0.933$, TZ $0.937$; MAP TZ $0.965$). Conversely, classifying on the five intermediate fractions alone---the restricted, free-water, and $D = 20\;\umsmm$ vascular pseudo-diffusion fractions removed---dropped NUTS detection to PZ $0.818$ / TZ $0.796$ (Table~\ref{tab:auc}), pinning down where the detection signal lives.

The individual-fraction AUCs (Table~\ref{tab:auc}) show the restricted and free-water bins to be the strongest single detectors (AUC $0.88$--$0.92$), but several intermediate fractions also discriminate on their own (up to $\sim$$0.85$), and even the four most poorly identified bins ($D = 0.5$--$1.5\;\umsmm$, posterior CV $> 0.7$) reach AUC $0.82$ (PZ) / $0.81$ (TZ) when used together. Which intermediate bin discriminates least flips between zones, reflecting the collinearity that leaves their weights poorly resolved within each ROI's posterior (Figure~\ref{fig:convergence}).
Detection is thus both carried and bounded by the restricted and free-water compartments, no intermediate fraction adding anything once those two are present.
Detection AUCs varied by $< 0.03$ across $C = 0.1$--$10$. Grade-classification AUCs are not reported---the $n = 29$ graded tumors are too few for a reliable estimate---and grade-dependent spectral changes are taken up in the Discussion.

\subsection{Agreement between ADC and the spectral discriminant}

Across ROIs, ADC was near-perfectly anti-correlated with the spectral discriminant in both zones and for both estimators (Figure~\ref{fig:adc_discriminant}), so ranking ROIs by ADC is almost identical to ranking them by the optimal spectral classifier. Using NUTS posterior-mean features at $C = 1.0$, the Spearman rank correlation---the measure relevant to detection, which depends only on how ROIs are ranked---was $\rho = -0.979$ (PZ, $n = 81$) and $-0.977$ (TZ, $n = 68$). The Pearson correlation, quantifying the linear near-identity of the scores themselves, was $-0.977$ ($95\%$ bootstrap CI $[-0.985,\,-0.964]$) and $-0.981$ ($[-0.988,\,-0.971]$). Constrained-MAP features gave equivalent values (Spearman PZ $-0.980$, TZ $-0.929$; Pearson $-0.972$, $-0.967$). The near-identity held across regularization strengths ($C = 0.1$ to $10$), indicating a property of the data rather than of any particular classifier setting.

Figure~\ref{fig:sensitivity} resolves the two readouts bin by bin.
In both zones the discriminant placed its two largest standardized coefficients on the restricted bin ($D = 0.25\;\umsmm$, normalized weight $+1.00$) and the free-water bin ($D = 3.0\;\umsmm$, $-0.63$ PZ and $-0.71$ TZ), and these were the only two whose $95\%$ bootstrap intervals excluded zero in both zones. Two intermediate bins additionally excluded zero in the transition zone ($D = 0.75$ and $1.0\;\umsmm$), with coefficients two to five times smaller, and none did so in the peripheral zone.
The ADC sensitivity vector $\partial\text{ADC}/\partial R_j$, evaluated at the mean tumor spectrum, was largest in magnitude and of opposite sign on the same two bins ($-1.74$ at $D = 0.25\;\umsmm$ and $+0.83$ at $D = 3.0\;\umsmm$, peripheral zone), but remained appreciable across the intermediate bins, where the discriminant does not.
Taken across the eight bins, the sensitivity vector and the discriminant coefficients were anti-correlated at $r = -0.79$ (PZ) and $-0.88$ (TZ), the negative sign following from tumor corresponding to \emph{low} ADC, so that the tumor-pointing sensitivity is $-\partial\text{ADC}/\partial\Rvec$. This correlation is computed over the eight bin weights and is a different quantity from the ROI-level score correlations above.
